\section{Introduction}
\label{sec:introduction}

Promotion planning is a central problem for businesses, such as retailers, online platforms, and service providers. Consumer goods companies in an industry study devoted nearly 20\% of their gross revenue to supporting retail promotions \citep{singerholston2017}. Better promotion decisions can also yield substantial revenue gains. In an evaluation using data from large retailers, \citet{cohenhillel2023} report average improvements ranging from 5.1\% to 15.6\% relative to existing practices. Promotion planning involves two key decisions. Frequency determines how often discounts are offered, while arrangement determines how the promotional periods are arranged in time. Businesses use both dispersed and concentrated patterns in practice. The movie theater chain Cinemark offers discounted tickets on Tuesdays at participating theaters \citep{cinemarkdiscounttuesdays}, while online retailer Amazon held its 2025 Prime Day event over four consecutive days \citep{amazon2025primeday}.

While the frequency of promotions has a clear impact on demand, arrangements of promotions can also substantially affect demand and revenue. Research shows that promotional patterns influence consumers' perceptions of deal frequency and that these perceptions are associated with willingness to pay \citep{krishna1991patterns}. Laboratory evidence also shows that purchase quantities vary across promotional patterns \citep{krishna1994behavior}. Consider a seller planning four promotion days in a month. Spreading them across four weeks exposes customers to a different price history than grouping them as a campaign at the beginning of the month. Also, different arrival times of customers will lead to different experiences of promotions. Under the concentrated arrangement, early arrivers may experience several consecutive promotions, while later arrivers may encounter none during the remainder of the month, by which time they may already have decided to drop out.

Furthermore, promotions within customer memory can affect future demand.
Empirical research documents persistent effects on price sensitivity and subsequent demand, with their direction depending on the customer and setting \citep{mela1997longterm,anderson2004,neslin2009dynamics}. Offering one additional promotion may have a diminishing or increasing effect on future demand. We hence note the following two channels. In the \textit{reference price channel}, past discounts lower the benchmark against which customers evaluate the current price, potentially reducing subsequent demand at a given price \citep{winer1986reference,mazumdar2005reference}. In the \textit{engagement channel}, promotions encourage customers to return, browse, or purchase again and may lead to higher demand in the future. A randomized experiment on the online retailer Alibaba documents increases in subsequent product views and purchase incidence \citep{dai2020alibaba}, while a supermarket field experiment finds that coupon redemption can reduce customer churn and increase the promoted product's customer lifetime value \citep{tsujikawa2025coupon}. Arrangement decisions therefore need to take these into account as well.

The central question of this paper is how a seller should choose both promotion frequency and arrangement to maximize long run average revenue when demand depends on the current price and remembered past promotions. Our analysis provides sufficient conditions that make both decisions solvable under finite customer memory.

This paper is related to several research areas. The literature on reference prices establishes how historical price perceptions influence promotion profitability and optimal price dynamics. \citet{greenleaf1995reference} studies profitable recurring promotions, and \citet{kopalle1996reference} derive conditions for constant or cyclic pricing under asymmetric responses to gains and losses. \citet{fibich2003reference} obtain explicit solutions for nonsmooth reference price models, while \citet{popescu2007} characterize pricing dynamics under general nonlinear reference effects. More recently, \citet{agrawal2024reference} establish near optimal markdown policies when reference prices reflect average past prices. These results depend on the reference formation mechanism and the response to perceived gains and losses. We allow for engagement effects as well as reference-price effects, and aim to characterize optimal arrangement and frequency when demand responds to the number of remembered promotions in complex ways.

There is also literature on applied promotion planning research which develops computational methods for optimizing promotion timing when demand depends on past prices. \citet{silvarisso1999} combine demand estimation, market simulation, and optimization to plan manufacturers' promotions. \citet{cohen2017promotion} develop a linear approximation that sums the profit increments from individual promotions, while \citet{cohen2021multiple} incorporate substitution and complementarity across items. \citet{baardman2019vehicles} study promotional vehicle scheduling under capacity constraints. These studies also establish exact solutions in specified cases, including sufficiently separated promotions in \citet{cohen2017promotion} and an unconstrained setting with two prices and substitutable items in \citet{cohen2021multiple}. Our analysis allows for general market share functions and derives sufficient conditions for optimal arrangement and frequency choices.

Another related body of literature establishes results on which features of past prices influence demand under customer memory. \citet{nasiry2011} characterize optimal pricing when the reference price combines the lowest and most recent prices. \citet{cohenhillel2023} establish conditions for optimal high and low pricing and provide an exact polynomial dynamic program for a single product with demand based on peak and recent prices. \citet{baek2025} characterize optimal cycles when customers remember the best recent offer and revenue satisfies a monotonicity condition. Our model allows for demand to depend on the number or counts of promotional periods encountered by a customer and characterizes the optimal arrangement and frequency. We also show how to handle the case where promotional effects depend on how recent the promotions are, without prescribing a particular functional form in the affine case in Section~\ref{sec:affine-location-dependent}.

Optimal sequencing literature provides the main mathematical precedent for our arrangement comparisons. \citet{hajek1985extremal} and \citet{altman2000multimodularity} establish optimality of regular sequences under multimodular costs, while \citet{altman2000balanced} also connect balanced routing to activity rate optimization. Their arguments link discrete convexity to the benefits of distributing repeated decisions evenly. However, their methods cannot be directly applied in our settings and we need to develop new techniques to characterize the optimal arrangements.

Technical difficulties of the promotion planning problem include the fact that each promotion enters several customer memory windows within one common price sequence. Also, nonlinear market share functions create interactions among promotions within each history, and transaction prices determine how those interactions contribute to revenue. A standard finite graph formulation establishes the existence of an optimal periodic policy, but its state space grows exponentially with memory depth, making computation expensive for long memories. Hence, it is an important and challenging question to characterize the optimal arrangements and frequency and find a strategy to solve the problem efficiently.

Our main contribution is a set of sufficient conditions for globally optimal promotion arrangements and frequencies. We use a finite telescoping representation of market share, with increments given by the marginal effect \(\Delta_{a,k}^j\) of increasing remembered promotions from \(k-1\) to \(k\), holding age and current price fixed. These marginal effects yield a decomposition of cycle revenue into a component fixed by cycle length and number of promotions and a sum of window contributions \(F_b\). Concave contributions favor balanced arrangements, in which the number of promotions differs by at most one across windows of the same length. Convex contributions favor a single promotion block. These conditions account for nonlinear demand, customer age, and transaction prices.

We can then solve the promotion frequency optimization problem under the same conditions. With concave window revenue contributions, revenue under balanced cycles is a concave, piecewise linear function of promotion frequency. An optimal frequency can be selected from an explicit finite set, yielding a policy that is globally optimal over all feasible infinite price paths and has a cycle of length at most maximum active age plus one. With convex window revenue contributions, a single block is optimal at a prescribed cycle length and number of promotions, while unrestricted frequency optimization yields a constant price policy.

We apply the general results to additively separable and multiplicative market share, two special cases with clearer economic interpretations. For additive demand, we distinguish engagement from reference price effects. Diminishing marginal engagement in exposure and age gives balanced optimality, with affine demand as a special case and a further extension allowing effects to vary across lags. Linear, loss neutral reference price effects give a constant price benchmark. Multiplicative demand models retention and churn through a Markov chain, with price and retention ratios determining arrangement and frequency regimes.

{In addition to the exact solutions, for the mixed-curvature scenarios, we provide two computational approximations with revenue guarantees. Finite searches over balanced and single block policies have loss bounds controlled by violations of the corresponding curvature conditions. Memory truncation reduces the graph and bounds loss through demand errors caused by omitted histories. }

{We examine implementation, sensitivity to demand assumptions, and computational approximation through three numerical experiments. We first illustrate how demand estimation and structural checks guide the selection of a pricing cycle. We then use reference-price demand to assess the consequences of ignoring the order of past prices. Finally, we evaluate balanced and single block policies and memory truncation alongside their revenue loss bounds to identify when these approximations are useful and how well the guarantees reflect policy performance.}

The remainder of the paper is organized as follows. Section~\ref{sec:model} presents the model and its periodic reformulation. Section~\ref{sec:general-decomposition} develops the revenue decomposition, arrangement conditions, and frequency optimization. Sections~\ref{sec:affine} and~\ref{sec:beyond-linear} examine additively separable and multiplicative market share functions. Section~\ref{sec:approximation} develops simple policy and memory truncation guarantees under mixed curvature. {Section~\ref{sec:usecase} examines implementation, sensitivity to price order, and approximation performance.} Section~\ref{sec:conclusion} concludes. The supplementary appendix contains additional proofs and {supplementary numerical results}.
